\documentclass[11pt]{article}
\usepackage[margin=1in]{geometry}
\usepackage{amsmath,amssymb}
\usepackage[hidelinks]{hyperref}
\hypersetup{pdftitle={The exact covering number C(70,21,2) = 15},
  pdfauthor={Troy Valiyaparampil and GPT-6.2 Astra}}
\setlength{\parindent}{0pt}
\setlength{\parskip}{5pt}
\title{Pair coverings of 70 points by blocks of size 21\\[3pt]
\large The exact covering number \(C(70,21,2)=15\)}
\author{Troy Valiyaparampil and GPT-6.2 Astra}
\date{25 November 2026}
\begin{document}
\maketitle
\begin{abstract}
We prove that in any pair covering by fourteen blocks, some block contains
more than three tenths of the points. A sparse weighting of the point
incidence patterns reduces the critical case to twelve four-element patterns
of equal weight and at most two five-element patterns. A coordinate outside
the latter patterns violates a parity condition. Previously published
fifteen-block coverings then give \(C(70,21,2)=15\). The accompanying
Lean~4 proof verifies the impossibility of a fourteen-block covering.
\end{abstract}

\section*{The weighted obstruction}
Let \([14]=\{1,\ldots,14\}\). A point of a covering determines a
\emph{pattern}, the set of indices of the blocks containing it. Covering
every pair means that distinct patterns intersect. Points with the same
pattern can be represented by one set whose weight is their combined mass.
Write \(d_i=\sum_{S\ni i} w_S\) for the load on index \(i\).

\textbf{Theorem.} For any pairwise-intersecting family of nonempty subsets
of \([14]\) with nonnegative weights summing to \(30\), some load exceeds \(9\).

\textit{Proof.} Suppose every load is at most \(9\). On this same family,
maximize total weight subject only to nonnegativity and the fourteen load
bounds. A maximum exists because the feasible region is closed and each
weight is at most \(9\). Choose a maximizer with as few positive weights as
possible. Its total is at least \(30\).

There are at most fourteen positive weights. Otherwise their incidence
vectors in \(\mathbb R^{14}\) have a nonzero linear dependence with
coefficient vector \(h\). Adding a sufficiently small multiple of \(h\) of
either sign preserves all loads and nonnegativity. Maximality implies
\(\sum_S h_S=0\). We can therefore continue in one direction until the
first positive weight becomes zero, preserving the maximum and contradicting
minimal support. Rescale the maximizing weights to total \(30\). All loads
are still at most \(9\). In what follows, discard the zero weights.

For a remaining pattern \(S\) of size \(r\), intersection with every other
pattern gives
\[
  9r \ \ge\ \sum_{i\in S} d_i
    =\sum_T |S\cap T|w_T
    \ \ge\ 30+(r-1)w_S .
    \tag{1}
\]
In particular every \(r\ge4\), each four-element pattern weighs at most
\(2\), and each five-element pattern weighs at most \(15/4\).
Summing the loads gives a second inequality,
\[
  \sum_S (|S|-4)w_S=\sum_{i=1}^{14}d_i-120\ \le\ 6 .
  \tag{2}
\]

Patterns of size at least five thus carry at most \(6\) units of weight.
At least \(24\) units lie on four-element patterns, so there are at least
twelve of these. Let \(n\) denote their number.

\emph{If \(n=14\)}, they account for the entire support and weigh at most
\(28\) altogether.

\emph{If \(n=13\)}, the four-element patterns weigh at most \(26\).
One further pattern must carry at least \(4\) units. A five-element pattern
weighs at most \(15/4<4\), while one of size at least six contributes at
least \(2\cdot4=8\) to (2).

\emph{If \(n=12\)}, all twelve four-element patterns must weigh exactly
\(2\), and the other patterns together weigh \(6\). Equality in (2) forces
every other positive pattern to have size five and every load to equal \(9\).
There are at most two five-element patterns. They cover at most ten of the
fourteen indices. At an index outside their union the load is twice the
number of four-element patterns through it. This even integer cannot equal
\(9\). The three cases exhaust the possibilities. \(\square\)

\section*{The covering number}
For integers \(v\ge k\ge2\), \(C(v,k,2)\) is the least number of
\(k\)-element subsets needed to cover every pair of a \(v\)-element set.
At \(v=70,\ k=21\), a point needs at least four blocks to meet its \(69\)
partners because a block supplies at most \(20\) partners. Counting
incidences gives the elementary lower bound
\(\lceil 70\cdot4/21\rceil=14\).

In a hypothetical fourteen-block cover give each point weight \(30/70\).
After combining repeated patterns, each load equals \(21(30/70)=9\),
contrary to the theorem. On the other hand, explicit fifteen-block
certificates are available from IRSI, Meliwat, and Arhan Canli (Quantavane)
[3--5]. Direct enumeration verifies every one of the \(\binom{70}{2}=2415\)
pairs. Hence \(\boxed{C(70,21,2)=15}\).

\section*{Relation to previous work}
Mills investigated coverings with few blocks [1]. Hall, Horsley and
Stinson [2] relate the covering ratio to the data limit \(L(m)\) and to
fractional matchings. The support reduction they attribute to F\"uredi
and their individual weight bound are the ingredients used above.
Their Theorem 1.5 at \(m=14\) gives
\((123+\sqrt{1297})/532<3/10\) as the comparison threshold.
The integer count of four-element patterns and the parity contradiction
exclude \(3/10\) itself. The minimum defining \(L(14)\) is attained [2],
so the theorem also gives \(L(14)>3/10\).

\section*{Machine verification of the lower bound}
Lean~4.24.0 checked \texttt{C7021.noFourteenBlockCover} in
\texttt{formal/Lemma.lean}. The theorem quantifies over every indexed family
of fourteen 21-element subsets of a 70-element set, allowing repeated blocks.
It proves that some pair of distinct points shares no block.

The formal proof starts each point at weight \(3/7\). When more than
fourteen points have positive weight, a rational linear dependence lets us
preserve every load and eliminate one positive weight without decreasing
the total. Rescaling gives the sparse weighting used in (1) and (2).
Lean then checks the rank counts and the parity contradiction.

\small
\clearpage
\section*{An optimal covering}
Each line below is a block. The labels run from 1 to 70. This certificate
is reproduced unchanged from Arhan Canli (Quantavane) [5], under
\href{https://creativecommons.org/licenses/by/4.0/}{CC BY 4.0}.
\begin{verbatim}
2 5 9 11 13 20 26 29 30 32 33 38 39 40 41 46 47 50 60 64 68
2 8 10 11 16 25 27 31 32 33 35 43 46 53 57 58 60 61 64 67 68
3 4 7 8 16 18 25 26 27 38 41 43 44 45 47 49 53 58 61 65 67
2 7 11 14 18 21 23 28 32 33 34 37 46 49 52 60 63 64 65 68 69
5 8 14 16 19 24 25 27 29 30 34 43 50 52 53 54 58 61 62 63 67
2 6 10 11 12 15 22 32 33 46 48 51 53 55 56 57 60 64 66 68 70
6 8 12 15 16 20 22 23 25 27 42 43 48 51 55 56 58 61 66 67 70
1 5 7 12 17 18 22 29 30 31 35 36 49 50 51 55 56 59 65 66 70
3 4 5 6 10 15 20 21 23 28 29 30 37 42 44 45 48 50 53 57 69
3 4 9 12 13 14 22 31 34 35 39 40 44 45 51 52 55 56 63 66 70
1 8 9 13 16 17 21 25 27 28 36 37 39 40 43 53 58 59 61 67 69
6 7 9 10 13 15 18 19 20 23 24 31 39 40 42 48 49 54 57 62 65
1 6 10 14 15 17 20 23 26 34 35 36 38 41 42 47 48 52 57 59 63
1 2 3 4 11 17 19 24 32 33 36 42 44 45 46 54 59 60 62 64 68
12 19 21 22 24 26 28 31 35 37 38 41 47 51 54 55 56 62 66 69 70
\end{verbatim}
Of the 2,415 pairs, exactly 2,107 occur once, 124 twice, 22 three times,
81 four times and 81 five times. Their multiplicities sum to
\(3150=15\binom{21}{2}\). The files from [3] and [4] give the same design
up to permutations of points and blocks. Explicit permutations and two
independent certificate verifiers accompany this note.

\textbf{Reproducing the Lean check.} Follow \texttt{formal/README.md}
to build the pinned project, then run \texttt{lake env lean Audit.lean}
from \texttt{formal/}. This checks the theorem statement and its axiom
dependencies. The fresh local replay in \texttt{formal/REPLAY.txt} reported
only \texttt{propext}, \texttt{Classical.choice} and \texttt{Quot.sound}.

\textbf{References and data.}
\begin{enumerate}
\item W.~H. Mills. Covering designs I. Coverings by a small number of
subsets. \emph{Ars Combinatoria} 8 (1979), 199--315.
\item J. Hall, D. Horsley and D.~R. Stinson.
\href{https://doi.org/10.1007/s10623-024-01418-6}{Bounds on data limits for
all-to-all comparison from combinatorial designs}.
\emph{Designs, Codes and Cryptography} 92 (2024), 2863--2879.
See Lemmas 3.1--3.4 and Section 5.
\item Institute for Recursive Self Improvement.
\href{https://irsi.ai/tables.html}{Table records and certificate archive},
September 2026.
\item Meliwat. \href{https://github.com/Meliwat/covering-hunt}{Covering hunt},
record \texttt{C70\_21\_2\_b15.txt}, October 2026.
\item A. Canli, Quantavane.
\href{https://quantavane.com/results/coverings/c-70-21-2/}{Covering design
with 15 blocks}, 22 October 2026.
\end{enumerate}
\textbf{Contributions.} Troy Valiyaparampil posed the research question
and requested formal verification. GPT-6.2 Astra, working in OpenAI Codex,
developed the lower-bound argument and drafted the paper. GPT-6.2 Sol
sessions reviewed the argument, checked certificate data and contributed
Lean proof modules. The fifteen-block constructions are credited in [3--5].
\end{document}
